\section*{الفصل \ref{ch:g7:sediments-and-fossils} --- الرواسب والأحافير}

\begin{solution}{exo:g7:sediments-and-fossils:1}
(أ) \omterm{def:g7:sediments-and-fossils:fossil}{القالب}؛ (ب) \omterm{def:g7:sediments-and-fossils:fossil}{الأحفورة الأثرية}؛ (ج) \omterm{def:g7:sediments-and-fossils:sediment}{الرواسب}؛ (د) \omterm{def:g7:sediments-and-fossils:fossil}{المصبوبة}؛ (ه) \omterm{def:g7:sediments-and-fossils:layer}{الطبقة
الصخرية}؛ (و) \omterm{def:g7:sediments-and-fossils:fossil}{الأحفورة}.
\end{solution}

\begin{solution}{exo:g7:sediments-and-fossils:2}
(ب) يموت الحيوان؛ (د) يغطي \omterm{def:g3:stones-pebbles-sand:sand}{الوحل} الصدفة بسرعة؛ (ج) تتراكم طبقات جديدة
وتتحول الصدفة إلى حجر؛ (أ) تبري \omterm{def:g5:water-shapes-the-land:erosion}{التعرية} الطبقات التي فوقها.
\end{solution}

\begin{solution}{exo:g7:sediments-and-fossils:3}
(أ) \omterm{def:g2:water-on-earth:ocean}{بحر} دافئ ضحل صافٍ. (ب) مياه ضحلة، على حافة \omterm{def:g1:landscapes-around-us:water}{بحيرة} أو \omterm{def:g1:landscapes-around-us:water}{نهر} أو \omterm{def:g2:water-on-earth:ocean}{بحر}. (ج)
غابة مستنقعية أو مستنقع. (د) يابسة تنمو عليها نباتات (وتطير فوقها
اليعاسيب)، قرب ماء ترسب فيه \omterm{def:g3:stones-pebbles-sand:sand}{الوحل}.
\end{solution}

\begin{solution}{exo:g7:sediments-and-fossils:4}
(أ) خطأ: معظم الكائنات الحية لا تترك أي \omterm{def:g7:sediments-and-fossils:fossil}{أحفورة}؛ فهي تتعفن أو تذوب قبل أن
تُدفن. (ب) صواب: إنه \omterm{def:g7:sediments-and-fossils:fossil}{أحفورة أثرية}. (ج) خطأ: أكبر الحبيبات وأثقلها تترسب
أولًا، \omterm{def:g3:stones-pebbles-sand:sand}{والوحل} الدقيق أخيرًا. (د) صواب.
\end{solution}

\begin{solution}{exo:g7:sediments-and-fossils:5}
(أ) الحجر الرملي \qty{120}{m}، والحجر الجيري \qty{180}{m}، والسجّيل
\qty{100}{m}. (ب) $120 + 180 + 100 = \qty{400}{m}$. (ج) $180 \div 400 = 0.45$:
\textbf{45~\%}. (د) لأن سطحه تآكل بفعل \omterm{def:g5:water-shapes-the-land:erosion}{التعرية}: فما بقي منه ليس الطبقة
كلها، ومقدار ما بقي يتغير من مكان إلى آخر.
\end{solution}

\begin{solution}{exo:g7:sediments-and-fossils:6}
(أ) من الأسفل إلى الأعلى: \omterm{def:g3:stones-pebbles-sand:pebble}{الحصباء}، ثم \omterm{def:g3:stones-pebbles-sand:sand}{الرمل}، ثم \omterm{def:g3:stones-pebbles-sand:sand}{الوحل} في الأعلى، تحت
الماء الصافي. (ب) حبيبات \omterm{def:g3:stones-pebbles-sand:sand}{الوحل} الدقيقة تغوص ببطء شديد: فهي لا تزال تسقط
بعد ساعة. (ج) مجموعة ثانية من الطبقات فوق الأولى: \omterm{def:g3:stones-pebbles-sand:sand}{رمل}، ثم \omterm{def:g3:stones-pebbles-sand:sand}{وحل} من جديد.
(د) \omterm{def:g3:stones-pebbles-sand:pebble}{الحصباء} التي في القاع تمامًا.
\end{solution}

\begin{solution}{exo:g7:sediments-and-fossils:7}
الأفضل: (ب)، صدفة صلبة، دُفنت بسرعة تحت \omterm{def:g3:stones-pebbles-sand:sand}{وحل} دقيق في ماء هادئ. ثم (د):
الصدفة صلبة ومدفونة، لكن \omterm{def:g5:water-shapes-the-land:flood}{الفيضان} قد يحركها أيضًا ويكسرها. والأقل: (أ)،
جسم طري بلا أجزاء صلبة، على \omterm{def:g3:stones-pebbles-sand:rock}{صخور} لا يترسب عليها \omterm{def:g3:stones-pebbles-sand:sand}{وحل} يغطيه. و(ج) ليست أفضل
بكثير: فالعظم يبقى مكشوفًا فيتحلل ويتآكل.
\end{solution}

\begin{solution}{exo:g7:sediments-and-fossils:8}
(أ) تجويف شكل الصدفة مطبوع على جدرانه: إنه \omterm{def:g7:sediments-and-fossils:fossil}{القالب}. (ب) \omterm{def:g7:sediments-and-fossils:fossil}{مصبوبة}، أي نسخة
حجرية من الصدفة. (ج) لا: فالصدفة ذابت؛ \omterm{def:g7:sediments-and-fossils:fossil}{والمصبوبة} مصنوعة من \omterm{def:g6:rocks-and-minerals:mineral}{معادن} جديدة.
(د) يمر الماء بسهولة بين حبيبات \omterm{def:g3:stones-pebbles-sand:sand}{الرمل} الكبيرة، فيستطيع أن يبلغ الصدفة
ويذيبها؛ أما في \omterm{def:g3:stones-pebbles-sand:sand}{الوحل} الدقيق فمروره أصعب بكثير.
\end{solution}

\begin{solution}{exo:g7:sediments-and-fossils:9}
(أ) $1.6 \times 1000 = \qty{1600}{m}$. (ب) $1600 \div 80 = 20$ طبقة.
(ج) $160 \div 1600 = 0.1$: \textbf{10~\%}.
\end{solution}

\begin{solution}{exo:g7:sediments-and-fossils:10}
(أ) لا: فالقدم لا تترك أثرًا إلا في أرض طرية، \omterm{def:g3:stones-pebbles-sand:sand}{رمل} مبلل أو \omterm{def:g3:stones-pebbles-sand:sand}{وحل}؛ والأرض
صارت \omterm{def:g3:stones-pebbles-sand:rock}{صخرًا} لاحقًا. (ب) تمحوها الموجة التالية أو \omterm{def:g5:water-shapes-the-land:flood}{الفيضان} التالي. (ج) كان
لا بد أن تغطيها \omterm{def:g7:sediments-and-fossils:sediment}{رواسب} جديدة وتملأها بسرعة، قبل أن تُمحى. (د) \omterm{def:g7:sediments-and-fossils:fossil}{أحفورة أثرية}:
فهو يُظهر ما فعله الحيوان، لا جزءًا من جسمه.
\end{solution}

\begin{solution}{exo:g7:sediments-and-fossils:11}
(أ) $20 + 5 = \qty{25}{mm}$. (ب) $1000 \div 25 = 40$ زوجًا: فالجزرة
تسجل \textbf{40 فيضانًا}. (ج) \omterm{def:g3:stones-pebbles-sand:sand}{الرمل} ثقيل: فهو يترسب بعد دخول \omterm{def:g1:landscapes-around-us:water}{النهر} إلى
\omterm{def:g1:landscapes-around-us:water}{البحيرة} بقليل، قرب \omterm{def:g5:water-shapes-the-land:source-mouth}{المصب}؛ ولا ينتقل منه بعيدًا إلا القليل. (د) بعيدًا عن
\omterm{def:g5:water-shapes-the-land:source-mouth}{المصب}، لا يبقى في الماء إلا \omterm{def:g3:stones-pebbles-sand:sand}{الوحل} الدقيق: فيترسب هناك ببطء، حين يكون
الماء هادئًا.
\end{solution}

\begin{solution}{exo:g7:sediments-and-fossils:12}
(أ) الأصداف صلبة: فهي تدوم زمنًا يكفي لتُدفن، وحيوانات كثيرة ذات أصداف
عاشت في قاع \omterm{def:g2:water-on-earth:ocean}{البحر}، حيث يترسب \omterm{def:g3:stones-pebbles-sand:sand}{الوحل} فوقها. (ب) الأجسام الطرية تتعفن قبل أن
تُدفن. (ج) لا: فالحيوانات الطرية لا تكاد تترك \omterm{def:g7:sediments-and-fossils:fossil}{أحافير} أبدًا، ولذلك فإن
القائمة المستمدة من \omterm{def:g7:sediments-and-fossils:fossil}{الأحافير} ناقصة؛ إنها تُظهر في الغالب الحيوانات ذات
الأصداف أو العظام.
\end{solution}

\begin{solution}{pb:g7:sediments-and-fossils:1}
\textbf{1.} الحجر الرملي، فالحجر الجيري، فالسجّيل، فالطباشير (الأول في
أسفل العمود).

\textbf{2.} \qty{120}{m} و\qty{180}{m} و\qty{100}{m}؛ والمجموع
\qty{400}{m}.

\textbf{3.} $180 \div 400 = 0.45$: \textbf{45~\%}.

\textbf{4.} لأن سطحه تآكل: فالجزء العلوي من الطبقة زال، ولا يستطيع أحد
أن يقيس ما لم يعد موجودًا.

\textbf{5.} في مياه ضحلة على حافة \omterm{def:g1:landscapes-around-us:water}{بحيرة} أو \omterm{def:g1:landscapes-around-us:water}{نهر} أو \omterm{def:g2:water-on-earth:ocean}{بحر}: فعلامات التموج
تتكون على \omterm{def:g3:stones-pebbles-sand:sand}{الرمل} في المياه الضحلة، وآثار الأقدام تُترك في \omterm{def:g3:stones-pebbles-sand:sand}{رمل} أو \omterm{def:g3:stones-pebbles-sand:sand}{وحل} طري
مبلل، كثيرًا ما يكون على شاطئ.

\textbf{6.} في \omterm{def:g2:water-on-earth:ocean}{بحر} دافئ ضحل صافٍ: الدليل هو المرجان.

\textbf{7.} هادئًا: \omterm{def:g3:stones-pebbles-sand:sand}{فالوحل} الدقيق لا يترسب إلا حيث يكاد الماء يسكن.

\textbf{8.} في البداية كان المكان شاطئًا، فيه مياه ضحلة \omterm{def:g3:stones-pebbles-sand:sand}{ورمل} طري مشت عليه
حيوانات. ثم غطاه \omterm{def:g2:water-on-earth:ocean}{البحر}: \omterm{def:g2:water-on-earth:ocean}{بحر} دافئ ضحل نما فيه المرجان وتراكمت الأصداف.
وبعد ذلك صار الماء أهدأ، وربما أعمق، فلم يعد يترسب إلا \omterm{def:g3:stones-pebbles-sand:sand}{الوحل} الدقيق، مع
أصداف الأمونيت، وهي حيوانات من أقارب الحبّار والأخطبوط.

\textbf{9.} $180 \times 100 = 18\,000$~سنتيمتر.

\textbf{10.} $120~\text{m} = 12\,000$~سنتيمتر، و$12\,000 \times 20 =
240\,000$~سنة.

\textbf{11.} $1~\text{m} = 100$~سنتيمتر، و$100 \times 100 = 10\,000$~سنة.

\textbf{12.} $180 \times 10\,000 = 1\,800\,000$: نحو
\textbf{1\,800\,000 سنة}، أي قرابة مليوني سنة. وهذا يعادل
$1\,800\,000 \div 240\,000 = 7.5$ مرة زمن الحجر الرملي.
\end{solution}
